\question{}

\begin{ابجد}
\item \textbf{درست}: در شبکه‌ی بیزی ساخته شده در الگوریتم \lr{Naive Bayes} تنها بین متغیر کلاس و متغیرهای ویژگی یال داریم و بین متغیرهای ویژگی یالی وجود ندارد.

\item \textbf{درست}: به عنوان مثال اگر در \lr{Spam/Ham Classification} در داده‌های آموزشی یک کلمه را در یک کلاس ندیده باشیم ولی در داده‌های تست باشد، مدل احتمال وقوع آن کلمه را صفر و احتمال کل صفر می‌شود. در \lr{smoothing} تصور می‌شود که هر کلمه را قبلا دیده‌ایم و این مشکل رخ نمی‌دهد.

\item \textbf{نادرست}: می‌توان از درخت‌های تصمیم برای مسائل \lr{Regression} نیز استفاده کرد. (که به این دسته‌ی خاص، \lr{Regression Trees} گفته می‌شود.) در این حالت متغیرها به جای اینکه عضو چند کلاس خاص باشند (مقدار گسسته) می‌توانند مقادیر پیوسته بگیرند.

\item \textbf{درست}: با افزایش عمق، درخت پیچیده‌تر شده و می‌تواند جزئیات بیشتری را یاد بگیرد. در این حالت درخت علاوه بر الگوهای کلی، نویزهای داده‌های آموزشی را نیز یاد می‌گیرد و می‌گوییم \lr{Overfitting} رخ داده.

\item \textbf{درست}: اگر یک ویژگی تعداد زیادی مقدار یکتا داشته باشد، با تقسیم کردن بر اساس آن ویژگی به گره‌های فرزند زیادی می‌رسیم که هر کدام تنها شامل چند نمونه هستند. پس آن گره‌ها \lr{pure} شده و آنتروپی کمی دارند. در این حالت آنتروپی شرطی نیز به صفر نزدیک شده و مقدار \lr{IG} افزایش می‌یابد. بنابراین درخت تمایل دارد آن ویژگی را انتخاب کند. (برای حل این مشکل، از معیار \lr{Gain Ratio} استفاده می‌شود.)

\item \textbf{نادرست}: درخت تصمیم می‌تواند فضای ویژگی‌ها را با مجموعه‌ای از برش‌ها به نواحی چندوجهی مختلف تقسیم کند و مرزهای غیرخطی را مدل کند.

\item \textbf{درست}: اگر مدل حتی روی داده‌های آموزشی هم عملکرد مناسبی نداشته باشد، یعنی مدل بیش از اندازه ساده بوده و نتوانسته الگوهای موجود در داده‌ها را به درستی یاد بگیرد. به این پدیده \lr{underfitting} گفته می‌شود.

\item \textbf{نادرست}: \lr{Hyper Parameter} ها از روی داده‌های \lr{Validation} یاد گرفته می‌شوند و اگر از داده‌های آموزشی برای این کار استفاده کنیم، \lr{overfitting} رخ می‌دهد. (می‌توانیم هم خودمان \lr{Hyper Parameter} ها را انتخاب کنیم.)

\end{ابجد}